\documentclass[12pt]{article}
%^ Start of document
\usepackage{times}  %Makes text TimesNewRoman
\usepackage{setspace} %Allows you to set single or double space 
\usepackage{biblatex} %Imports biblatex package
\usepackage{graphicx} % Required for inserting images
\usepackage{authblk}
\usepackage[colorlinks=true, urlcolor=blue, linkcolor=blue]{hyperref}
\usepackage{subfig}
\usepackage{float}
\usepackage{tabularx}
\usepackage{multirow}
\usepackage[paperheight=11in,
   paperwidth=8.5in,
   top=20mm,
   bottom=20mm,
   left=40mm,
   right=40mm]{geometry}
\usepackage{moresize}
\usepackage{tikz}
\usepackage{xcolor}
\usepackage{physics}
\usepackage{amssymb}
\usepackage{mathrsfs}
\usepackage{amsmath}
\usepackage{ragged2e}
\usepackage{comment}
\usepackage{xcolor}
\addbibresource{seminar.bib}
\begin{document}
\singlespacing  %Sets single spacing for whole document
\title{Globular Clusters - Reddell}

\author[1]{Zachary Tullier}
\affil[1]{Student\\Physics Department\\ University of Houston - Clear Lake \\Houston, Texas\\U.S}
\date{}
\maketitle




\section{Introduction}
    \textbf{Speaker:} Dr. Brandon Reddell, an astrophysicist at the University of Houston-Clear Lake, delivered a detailed lecture on the role of globular clusters in estimating the lower limit of the universe's age. His expertise bridges theoretical astrophysics with observational techniques, including methods accessible to amateur astronomers. This lecture focused on how globular clusters, some of the oldest objects in the universe, serve as natural clocks to estimate the universe's age. Emphasized using Hertzsprung-Russell (HR) diagrams and stellar isochrones to approximate cluster ages.


\section{The Hertzsprung-Russell (HR) Diagram}
    The HR diagram is a fundamental tool in astronomy, plotting stars by luminosity and temperature (or spectral class). It provides a snapshot of a star cluster's population, showing stages of stellar evolution. The vertical axis is luminosity (intrinsic brightness), often expressed as absolute magnitude. The horizontal axis is temperature or spectral class, with temperature decreasing from left to right. Apparent Magnitude (\( m \)) is a star's brightness as seen from Earth. Absolute Magnitude (\( M \)) is a star's brightness if placed at 10 parsecs. Distance Modulus is \( m - M = 5 \log(d) - 5 \), where \( d \) is the distance in parsecs. As indicated in the diagram, main sequence stars burn hydrogen in their cores, giants and supergiants have exhausted core hydrogen, and white dwarfs are compact remnants of low-mass stars.
    
\section{Globular Clusters: Unique Characteristics}
    Globular clusters are dense, spherical collections of stars bound by gravity, located in the halo and plane of galaxies and the contain tens of thousands to millions of stars. Stars in a globular cluster formed around the same time and from the same material. They have differences in appearance that arise from evolutionary stages and mass. Globular clusters are among the universe's oldest objects, with ages exceeding 12 billion years. The can provide a reliable lower limit for the age of the universe. The Turnoff Point is the point on the HR diagram where stars leave the main sequence, indicating core hydrogen exhaustion. The location of the turnoff point corresponds to the cluster's age: more massive stars evolve off the main sequence earlier.


\textbf{Determining Cluster Ages:}
    Observe the cluster in B (blue) and V (visual) filters to measure apparent magnitudes (\( m_B \) and \( m_V \)) then calculate the color index (\( B-V \)) for each star. Plot the HR diagram (magnitude vs. \( B-V \)) and correct for extinction then compare the observed HR diagram with theoretical stellar isochrones to estimate the cluster's age. Images are taken through B and V filters to capture magnitudes and color indices. Calibrations are performed against known reference stars to ensure accuracy.
 
\section{Case Study: Messier 12 (M12)}
    Observations of M12 show that it is ~15,600 light-years from Earth. Magnitudes and color indices derived from calibrated imaging. Turnoff point analysis yielded an estimated age of ~12.3 billion years for M12.
    Distance modulus (\( m - M = 13.4 \)) and bolometric corrections refined luminosity values. The estimated age aligns closely with published values (See below), demonstrating the reliability of the methodology.
    \newline Age Estimates for Other Globular Clusters
    \begin{itemize}
        \item M13: 11.6 billion years.
        \item M2: 13 billion years.
        \item M15: 12 billion years.
        \item HP 1: 12.8 billion years.
        \item M22: 12 billion years.
    \end{itemize}
    The oldest clusters consistently indicate a formation ~13 billion years ago, shortly after the Big Bang. These methods, like B and V filter imaging and HR diagram plotting, are accessible to skilled amateurs. Globular cluster ages complement cosmic microwave background measurements and galaxy redshift surveys, providing independent benchmarks for the universe's age.


\section{Challenges and Future Directions}
    Correcting for extinction and atmospheric effects is crucial for accuracy and very difficult to accomplish. Space-based telescopes like the James Webb Space Telescope (JWST) promise more precise cluster age estimates. Combining globular cluster data with other cosmological tools will refine our understanding of the early universe.


\section{Summary and Takeaways}
Observing globular clusters through B and V filters, correcting for extinction and atmospheric distortions, and using HR diagrams and stellar isochrones to estimate cluster ages is possible for amateur astronomers. Established calculations date the oldest clusters consistently range from 12 to 13 billion years, providing a lower limit on the universe's age. Messier 12 (M12) serves as a case study with an age of ~12.3 billion years.

\end{document}
